\scanpage{47}
\mspage{2}

Condition pour que $N'=\uline{\operatorname{Cent}}G$

\unsure{a|*}) Il faut que $N'\subset\uline{\operatorname{Cent}}G$, i.e. op. \inserted{$\theta'$} de $M'$ sur $N'$
triviale ($\Longrightarrow N'$ commutatif et $d'(N')\subset\uline{\operatorname{Cent}}(M')$)

b) D'ailleurs, l'image du centre de $G$ dans $K$ \inserted{$=G/\pi_1$} est
contenue dans $\uline{\operatorname{Cent}}(K)$ \inserted{$\supset N'/\pi_1$ i.e. son image dans $M$ contenue dans $\uline{\operatorname{Cent}}(K)/(N'/\pi_1)$}, et aussi dans
$\operatorname{Ker}(\theta:M\longrightarrow\uline{\operatorname{Aut}}_{\mathrm{gr}}(N))$. Soit
\[
\zeta=\deleted{\uline{\operatorname{Cent}}(M)}\inserted{\uline{\operatorname{Cent}}(K)/(N'/\pi_1)}\cap\operatorname{Ker}\theta
\]
et soit $\overline{N}'$ \deleted{\unsure{son}} l'image inverse dans $G$ \inserted{le ss-groupe invariant de $G$ \unsure{déf.}},
\inserted{$(N'\subset\uline{\operatorname{Cent}}(G)\subset\overline{N}')$}
qui est donc
une extension de $\zeta$ par $N'$. L'opération de
$G$ sur $\overline{N}'$, via $M$, par automorphismes, est triviale
sur $\unsure{N'}$ et sur $\zeta=\overline{N}'/N'$ \inserted{(car $\zeta\subset\uline{\operatorname{Cent}}(M)$)}, donc se définit par un
hom
\[
M\longrightarrow\uline{\operatorname{Hom}}_{\mathrm{gr}}(\zeta,N')
\]
prenant en fait ses valeurs dans $\uline{\operatorname{Hom}}_{\mathrm{gr}}(\zeta,\pi_1)$
\deleted{$\unsure{\operatorname{Hom}_{\mathrm{gr}}(M,\pi_1)}$}
(lequel, par construction de $\zeta$, est aussi trivial
sur $d(N)\subset M$ ($\zeta\subset\operatorname{Ker}\theta$). Donc l'hom
se factorise en
\[
\varphi_G:\pi_0\longrightarrow\uline{\operatorname{Hom}}_{\mathrm{gr}}(\zeta,\pi_1)
\qquad\deleted{\unsure{Z^1}}
\]
qui s'interprète \inserted{aussi} comme
\[
\psi_G:\zeta\longrightarrow\uline{\operatorname{Hom}}_{\mathrm{gr}}(\pi_0,\pi_1)
\qquad\deleted{\unsure{\operatorname{Hom}_{\mathrm{gr}}}}
\]
Soit $\zeta_G=\operatorname{Ker}\psi_G$, et $\overline{N}'_G$ l'image inverse de
$\zeta_G$ dans $G$, il est clair alors que
\[
\overline{N}'_G=\uline{\operatorname{Cent}}(G)
\]
donc $N'=\uline{\operatorname{Cent}}(G)$ équivaut à la condition
\[
\boxed{\zeta_G=\operatorname{Ker}\psi_G=1\quad\text{i.e. }\psi_G\text{ injectif}}
\]

Soit $G_1$ une deuxième sol. du Pb initial,
différant de $G$ par une extension \deleted{$\unsure{E}$} $E$ de $\pi_0$ par $\pi_1$
(op. triviale de $\pi_0$ sur $\pi_1$ i.e. extension centrale),

\margin{\uline{\unsure{donc plus}} :
\diagram{\begin{tikzcd}[ampersand replacement=\&,column sep=small,row sep=small]
\zeta=\operatorname{Ker}\theta\cap\operatorname{Cent}M\arrow[d]\\
Z^1(\pi_0,\pi_1)
\end{tikzcd}}
$M\simeq G/N'$\par
$\zeta\subset\unsure{\zeta'}=\operatorname{Im}\uline{\operatorname{Cent}}(K)$}
\mspage{46}
